\documentclass[11pt]{article}
\usepackage{icmlko}
\begin{document}

\title{복도를 다 세어 보니 만화가 제일 심하지 않았다}
\author{월드모델 실험실}
\date{노트 363 · 2026년 8월 1일}
\maketitle

\begin{abstract}
노트 362가 만화의 해를 ``복도''로 설명했다 --- 1,783행이 공유 축 한쪽 끝에
쌓여 있어 거기서 가르는 분기가 만들어지고 다른 도메인 행도 같이 내려간다.
\textbf{그 설명이 맞다면 복도가 더 심한 도메인은 더 해로워야 한다.} 그래서
(도메인 $\times$ 축) 전부에 같은 것을 쟀다 --- \emph{복도 위험} $=$ 학습 중
몫 $\times$ 치우침(관측 행 중앙값 차). \textbf{예상이 깨졌다.} 상위는
세계애니 \texttt{venue\_prominence}(0.1922)와 모바일
\texttt{entry\_friction}(0.1486)이고 \textbf{만화는 11위}(\texttt{venue\_
prominence} 0.0539)다. \texttt{entry\_friction} 의 만화는 상위 18에 들지도
않는다. \textbf{판을 올리는 것은 만화 빼기인데, 복도가 제일 심한 것은
만화가 아니다.} 노트 362의 표현도 정정한다 --- 그 글이 ``만화가 한쪽 끝에
통째로 쌓여 있다''고 쓸 때 견준 ``나머지''에 \emph{관측 안 된 행의 중립값
0.5} 가 섞여 있었다. 관측 행끼리만 견주면 세계애니도 중앙 0.00 으로
똑같이 끝에 있다(그 글의 표에는 이미 그렇게 적혀 있다). 남는 가설은
회계다 --- \textbf{만화만 한 번도 채점이 안 된다.} 세계애니·모바일도 복도를
만들지만 제 점수(가중 11\% · 16\%)로 값을 치르고, 만화는 아무 값도 안 낸다.
공정한 검사를 걸어 뒀다.
\end{abstract}

\section{복도 위험을 다 센다}

노트 362가 만화 하나에서 본 것을 눈금으로 만들었다. (도메인, 축)마다
\[ \text{복도 위험} = \frac{\text{그 도메인 학습 행}}{\text{전체 학습 행}}
   \times \big|\,\mathrm{med}(\text{그 도메인 관측값}) - \mathrm{med}(\text{나머지 관측값})\,\big| \]
분기의 이득이 대체로 무리 크기에 비례하고, 무리가 갈라지려면 값이 떨어져
있어야 한다는 두 조건을 곱한 것이다. \textbf{관측된 행끼리만 견준다} ---
중립값 0.5 를 섞으면 관측률이 낮은 축이 저절로 ``치우친'' 것처럼 보인다.

\begin{center}\includegraphics[width=\linewidth]{figs/map.pdf}\end{center}

\begin{center}
\begin{tabular}{rllrrr}
\hline
\# & 축 & 도메인 & 몫 & 치우침 & 위험 \\
\hline
1 & \texttt{venue\_prominence} & 세계애니 & 25.2\% & 0.76 & \textbf{0.1922} \\
2 & \texttt{entry\_friction}   & 모바일   & 14.9\% & 1.00 & \textbf{0.1486} \\
3 & \texttt{entry\_friction}   & 세계애니 & 25.2\% & 0.37 & 0.0940 \\
4 & \texttt{venue\_prominence} & 웹툰     & 20.1\% & 0.40 & 0.0795 \\
5 & \texttt{media\_push}       & 애니     & 14.0\% & 0.55 & 0.0769 \\
\multicolumn{6}{c}{$\cdots$} \\
\textbf{11} & \texttt{venue\_prominence} & \textbf{만화} & 17.0\% & 0.32 & \textbf{0.0539} \\
13 & \texttt{goods\_scale}     & 만화     & 17.0\% & 0.30 & 0.0519 \\
\hline
\end{tabular}
\end{center}

\textbf{만화는 11위이고, 노트 362가 주인공으로 쓴 \texttt{entry\_friction}
은 상위 18에 없다.}

\section{노트 362의 표현을 정정한다}

노트 362는 만화 행과 ``나머지''를 68열에서 비교해 셋만 떨어져 있다고
적었는데, 그 ``나머지''는 \emph{설계행렬}이라 축이 없는 도메인의 중립값
0.5 가 그대로 들어 있었다. 그래서 만화(관측 100\%, 값 0.00)가 실제보다
멀리 있어 보였다.

\textbf{그 글의 도메인별 표는 옳았다} --- 거기엔 세계애니도
\texttt{entry\_friction} 중앙 0.00 으로 적혀 있고, 본문도 ``세계애니는
중앙이 똑같이 0.00 이고 2,648행이다''라고 썼다. \textbf{틀린 것은 그림과
초록의 강조}였다. 관측 행끼리 재면 만화는 그 축에서 특별하지 않다.

\textbf{규약: 도메인 사이 값을 견줄 때는 설계행렬이 아니라 관측 행으로
견준다.} 중립 대입은 모형에 주는 것이지 비교에 쓰라고 있는 값이 아니다.
노트 266 · 354 · 355가 ``결측이 버는가''를 물었던 것과 같은 집안이고,
이번엔 그 결측이 \emph{비교}를 오염시켰다.

\section{그러면 왜 만화인가}

복도 위험이 해를 그대로 예측하지 않는다. 남는 차이는 하나다.

\begin{center}
\begin{tabular}{lrrrl}
\hline
도메인 & 학습 & 유보 & 판 가중 & 채점되나 \\
\hline
세계애니 & 2,648 & 300 & 11.0\% & 된다 \\
모바일   & 1,559 & 441 & 16.2\% & 된다 \\
웹툰     & 2,106 & 711 & 26.1\% & 된다 \\
\textbf{만화} & \textbf{1,783} & \textbf{6} & \textbf{0\%} & \textbf{안 된다} \\
\hline
\end{tabular}
\end{center}

\textbf{만화만 값을 안 낸다.} 세계애니가 복도를 만들어 남에게 해를 끼쳐도
제 300행을 잘 맞혀 판의 11\%를 벌어 온다. 만화는 같은 일을 하고 아무것도
안 벌어 온다 --- 유보 6행이라 \texttt{evaluate} 문턱 20을 못 넘는다.
\textbf{해가 아니라 회계다.}

이 가설은 ``만화가 특별히 나쁘다''가 아니라 ``만화의 몫이 순이익에서
빠져 있다''고 말한다. 그러면 세계애니 학습을 빼도 \emph{남의 점수}는
비슷하게 오를 것이고, 다만 세계애니 자신의 점수가 무너져 판 전체로는
손해일 것이다.

\textbf{공정한 검사를 걸었다.} 세 팔(현행 · 만화 학습 뺌 · 세계애니 학습
뺌)을 \emph{만화와 세계애니를 뺀 아홉 도메인}으로만 채점한다. 대상과
가중이 세 팔에서 같으므로 ``남에게 준 해''만 남는다. 미리 적었다 ---
\textbf{복도 위험이 해를 예측하면 세계애니 쪽이 더 커야 한다}(0.1922 대
0.0539). 돌고 있다.

\section{남은 것}

\begin{center}
\begin{tabular}{ll}
\hline
갈래 & 지금 아는 것 \\
\hline
\textbf{남에게 준 해} & 아홉 도메인 판으로 만화 대 세계애니. \textbf{돌고 있다} \\
문턱이 정확히 어디인가 & 1.7\%\textasciitilde{}8.5\% 사이(노트 362). 4\% 팔 하나면 좁아진다 \\
팝업 실측 41행 & 노트 360. 결정 셋을 다 팝업 65행이 내렸다 \\
만화 2026 스무 건 & 노트 356. 채점되게 만들면 회계가 닫힌다 \\
\hline
\end{tabular}
\end{center}

마지막 줄이 이 노트에서 새 뜻을 얻었다. 만화 2026 스무 건은 ``열한 번째
도메인''을 얻는 일이기만 한 것이 아니라 \textbf{이미 치르고 있는 값을
회수하는} 일이다.

\end{document}
